% Generated draft: requires mathematical and source review.
\documentclass[11pt]{article}
\usepackage[margin=25mm]{geometry}
\usepackage{fontspec}
\setmainfont{Latin Modern Roman}
\usepackage{amsmath,amssymb,mathtools,unicode-math}
\setmathfont{Latin Modern Math}
\usepackage{enumitem,needspace,xurl,hyperref}
\setlength{\parindent}{0pt}
\setlength{\parskip}{5pt}
\setlength{\emergencystretch}{3em}
\newcommand{\N}{\mathbb{N}}
\newcommand{\Z}{\mathbb{Z}}
\newcommand{\Q}{\mathbb{Q}}
\newcommand{\R}{\mathbb{R}}
\newcommand{\C}{\mathbb{C}}
\newcommand{\F}{\mathbb{F}}
\newcommand{\A}{\mathbb{A}}
\newcommand{\PP}{\mathbb{P}}
\newcommand{\E}{\mathbb{E}}
\newcommand{\eps}{\varepsilon}
\newcommand{\dd}{\,\mathrm{d}}
\begin{document}
\section*{UnsolvedMath ID: 20000465}
\textbf{Source problem code:} catalogue entry
\subsection{Definitions and mathematical statement}
Question 32, attributed to Voloch, concerns several different height-ordered families. For any fixed number field \(K\) and finite height family \(\mathcal F_H\) of hypersurfaces, put
\[
L(H)=\#\{X\in\mathcal F_H:X(K_v)\ne\varnothing\text{ for every place }v\},\quad
R(H)=\#\{X\in\mathcal F_H:X(K)\ne\varnothing\},
\quad B(H)=L(H)-R(H).
\]
A hypersurface satisfies the Hasse principle if everywhere local solubility implies a \(K\)-point. Hence its failure count is \(B(H)\); the unconditional proportion satisfying the principle is \(1-B(H)/\#\mathcal F_H\), whereas the rationally soluble proportion is \(R(H)/\#\mathcal F_H\), and the conditionally soluble proportion is \(R(H)/L(H)\) when \(L(H)>0\). These quantities must not be identified.
\begin{enumerate}
\item For all degree-\(d\) hypersurfaces in \(\PP^n_K\), determine the limiting proportions as \(H\to\infty\), and the degree/dimension circumstances in which failures have density zero. Coefficients up to scalar form \(\PP^{N-1}(K)\), \(N=\binom{n+d}{d}\); a precise height and treatment of representatives must be chosen. Over \(\Q\), one possible convention is primitive integral coefficient vectors up to overall sign, with maximum coefficient at most \(H\). The source allows general fixed \(K\) and does not supply its height normalization. Colliot-Th\'el\`ene explicitly questions whether almost-all satisfaction, rather than the opposite behavior, should be expected.
\item For diagonal forms over \(\Q\), consider coefficient vectors \(\mathbf a\in\Z^s\setminus\{0\}\), \(|a_i|\leq A\), and
\[
X_{\mathbf a}:\quad a_1x_1^d+\cdots+a_sx_s^d=0\quad\text{in }\PP^{s-1}.
\]
The tentative suggestion of Wooley and Venkatesh is an asymptotic involving a product of local densities, at least for \(s\geq3d+1\). With probability-normalized coefficient measures, the natural rational-solubility version asks whether
\[
\lim_{A\to\infty}\frac{\#\{\mathbf a:|a_i|\leq A,\ \mathbf a\ne0,\ X_{\mathbf a}(\Q)\ne\varnothing\}}
{(2\lfloor A\rfloor+1)^s-1}
=\delta_\infty\prod_p\delta_p.
\]
Here \(\delta_p\) is the Haar probability of \(\Q_p\)-solubility for coefficients in \(\Z_p^s\), and \(\delta_\infty\) uses uniform coefficient measure on \([-1,1]^s\). Existence/convergence and normalization of the product are part of what must be specified. The printed expression instead uses \(A^{-s}\) times the count and says ``satisfies HP'': it has a different box-volume normalization and, taken literally, counts locally insoluble forms as satisfying HP. Thus the displayed precise variant must not be declared a verbatim equivalent of that tentative expression. The source contrasts this proposal with historical individual-form circle-method results in a range written \(s\gg d\log d\).
\item For plane cubics \((n,d)=(2,3)\), determine the rational-solubility proportion conditional on everywhere local solubility and any decay rate if it vanishes. Stoll proposes that the conditional proportion is zero; Mazur asks for an exponent, and Stoll requests experimental evidence. De Jong asks what such conclusions imply about \(3\)-Selmer groups of elliptic curves. Translating counts of plane models into Selmer statistics requires specified Jacobians, heights, and model multiplicities; no such translation is given in the source.
\end{enumerate}

The attached aid reports different later outcomes for these families: a density equality for full Fano hypersurface families over \(\Q\), except cubic surfaces; a diagonal asymptotic only in the overlap of the Bright--Browning--Loughran local-density and Br\"udern--Dietmann exceptional-set ranges; and Bhargava's positive-proportion result contradicting Stoll's proposed zero conditional density for plane cubics. Complete theorem statements and the overlap range are not supplied, so these are reported literature claims, not blanket assertions for all \(K,n,d,s\).

The aid also supplies a height-transfer claim: if \(\mathcal T_H\subseteq\mathcal R_{\beta H}\), the reference-family failure count is at most \(Ct^m e(t)\), and the target-family locally soluble count is at least \(cH^m\), where \(C,c,\beta>0\) are independent of \(H\), then the target conditional failure proportion is at most
\[
\frac{C\beta^m}{c}\,e(\beta H).
\]
It reports an application from a Euclidean coefficient norm to maximum coefficient height giving \(O((\log H)^{-1/(48n)})\), using a Browning--Le Boudec--Sawin bound. The application requires positive local mass and appropriate size of the target family; it does not justify restricting an ambient density theorem to arbitrary thin or zero-local-mass subfamilies. Its complete geometric range and analytic hypotheses are not provided.
\subsection{Short English statement}
Determine densities of rational points and Hasse-principle failures in full hypersurface families, diagonal families, and plane cubics. The draft separates unconditional satisfaction from solubility conditional on local points, preserves the tentative historical claims, and records the supplied later qualifications and height-transfer bound.
\subsection{Sources}
\begin{itemize}
\item Voloch, Question 32, AIM workshop \emph{Rational and integral points on higher dimensional varieties}: \url{https://aimath.org/WWN/qptsurface2/qptsurface2.pdf}. The discussion names Colliot-Th\'el\`ene, Poonen, Wooley, Venkatesh, de Jong, Stoll, and Mazur.
\item Local registry \texttt{data/unsolvedmath/problems.json}, permanent ID \(20000465\), exact source code \texttt{catalogue entry}; canonical \texttt{aim-algebraic-number-theory-notes.json}, notes[116]. The supplied verification identifies the damaged item number as 32.
\item The attached aid names Bright--Browning--Loughran, Br\"udern--Dietmann, Bhargava, and Browning--Le Boudec--Sawin, without full bibliographic citations or theorem hypotheses. Attempt 1 is explicitly machine-generated, not peer reviewed, and requires independent verification before citation.
\end{itemize}
\end{document}
